%==============================================================================
% GLOSSARY DEFINITIONS
%==============================================================================

\newglossaryentry{ai_definition}{
  name={Artificial Intelligence},
  description={The simulation of human cognitive processes---including reasoning,
    learning, and self-correction---by computer systems. In the context of \nabta,
    AI encompasses the deep learning models and language model components that
    drive automated plant disease diagnosis and consultation}
}

\newglossaryentry{deeplearning}{
  name={Deep Learning},
  description={A subfield of machine learning based on artificial neural networks
    with multiple representation-learning layers. Deep learning is the core
    enabling technology for the convolutional and segmentation models deployed
    in the \nabta pipeline}
}

\newglossaryentry{cnn_def}{
  name={Convolutional Neural Network},
  description={A class of deep learning architecture specifically designed for
    processing grid-structured data such as images. CNNs employ learnable spatial
    filters to extract hierarchical feature representations from input images}
}

\newglossaryentry{segmentation}{
  name={Semantic Segmentation},
  description={A computer vision task in which every pixel of an input image is
    assigned a class label. In \nabta, semantic segmentation is performed by the
    U-Net model to isolate the leaf foreground from background content}
}

\newglossaryentry{objectdetection}{
  name={Object Detection},
  description={A computer vision task that simultaneously localises and classifies
    objects within an image, typically by producing bounding boxes accompanied by
    class labels and confidence scores}
}

\newglossaryentry{explainableai}{
  name={Explainable Artificial Intelligence},
  description={A collection of methods and techniques that enable the outputs
    of AI systems to be understood and interpreted by human observers. \nabta
    employs Grad-CAM as its primary explainability mechanism}
}

\newglossaryentry{transferlearning}{
  name={Transfer Learning},
  description={A machine learning technique in which a model pre-trained on a
    large dataset is adapted to a related task with a smaller domain-specific
    dataset. The U-Net in \nabta uses an EfficientNet-B0 backbone pre-trained
    on ImageNet}
}

\newglossaryentry{vectordb}{
  name={Vector Database},
  description={A specialised database designed to store and retrieve large
    numerical vectors efficiently, typically using approximate nearest-neighbour
    algorithms. \nabta uses a vector database to index knowledge base embeddings
    for the RAG consultation module}
}

\newglossaryentry{pathogen}{
  name={Pathogen},
  description={A biological agent capable of causing disease in a host organism.
    In an agricultural context, plant pathogens include fungi, bacteria, viruses,
    and parasitic organisms responsible for crop diseases}
}

\newglossaryentry{roi_def}{
  name={Region of Interest},
  description={A spatially bounded sub-region of an image selected for further
    analysis. In the \nabta pipeline, the YOLOv8 detection model identifies the
    infected leaf region as the region of interest for downstream processing}
}

\newglossaryentry{augmentation_def}{
  name={Data Augmentation},
  description={A set of image transformation techniques applied during model
    training to artificially increase dataset diversity and reduce overfitting.
    \nabta applies rotations, flips, colour jitter, and random erasing among
    other augmentation operations}
}

\newglossaryentry{inference}{
  name={Inference},
  description={The process of applying a trained machine learning model to new
    input data to generate predictions. In \nabta, inference refers to the
    sequential execution of the YOLOv8, U-Net, and CNN models on a submitted
    leaf image}
}

\newglossaryentry{precision}{
  name={Precision},
  description={A classification metric defined as the ratio of true positive
    predictions to all positive predictions made by the model:
    $\text{Precision} = \text{TP} / (\text{TP} + \text{FP})$}
}

\newglossaryentry{recall_def}{
  name={Recall},
  description={A classification metric defined as the ratio of true positive
    predictions to all actual positive instances in the dataset:
    $\text{Recall} = \text{TP} / (\text{TP} + \text{FN})$}
}

\newglossaryentry{sus_def}{
  name={System Usability Scale},
  description={A standardised ten-item questionnaire used to assess the perceived
    usability of a software system. Scores range from 0 to 100, with scores above
    80.3 classified as Excellent}
}
